\input _template

\setlanguage{EN}

\Title{Inverses Modulo $p$}

\MathcompsLink{inverses-mod-p}

\Author{Patrik Bak}

\sec Introduction

We assume familiarity with congruences modulo a prime~$p$. In this handout we build a small theory of multiplicative inverses modulo~$p$ and use it to solve problems.

\sec Theory

Throughout, $p$ denotes a prime.

\secc Recap

We freely use the following standard facts about congruences modulo~$p$:

\begitems
\i Congruences may be added, subtracted, and multiplied: if $a \equiv b$ and $c \equiv d \pmod{p}$, then $a \pm c \equiv b \pm d$ and $ac \equiv bd \pmod{p}$.
\i Cancellation: if $a \not\equiv 0 \pmod{p}$ and $ax \equiv ay \pmod{p}$, then $x \equiv y \pmod{p}$.
\i For any $a \not\equiv 0 \pmod{p}$, the residues $a, 2a, \ldots, (p-1)a$ are a permutation of $1, 2, \ldots, p-1$ modulo~$p$.
\i Fermat's Little Theorem: for any $a \not\equiv 0 \pmod{p}$, we have $a^{p-1} \equiv 1 \pmod{p}$.
\i Order: the order of $a \not\equiv 0 \pmod{p}$ is the smallest positive integer $d$ with $a^d \equiv 1 \pmod{p}$. It divides $p - 1$.
\i Primitive root: there exists $g \not\equiv 0 \pmod{p}$ of order $p - 1$. Every $a \not\equiv 0 \pmod{p}$ equals $g^k$ for some $k$.
\enditems

\secc Inverses

\EnvId{YsoLc1b4Fqes4Dd5A8svC-existence-and-uniqueness-of-inverse}
\Theorem{Existence and Uniqueness of Inverse}{
    Let $a \not\equiv 0 \pmod{p}$. Then there exists a unique residue $a^{-1} \pmod{p}$ with
    $$a \cdot a^{-1} \equiv 1 \pmod{p}.$$
    We write $\frac{1}{a}$ for $a^{-1}$, and $\frac{b}{a}$ for $b \cdot a^{-1}$.
}{
    Since $a \not\equiv 0 \pmod{p}$, the residues $a, 2a, \ldots, (p-1)a$ are a permutation of $1, 2, \ldots, p-1$ modulo $p$. Hence exactly one of them is congruent to $1$, which gives a residue $a^{-1}$ with
    $$
    a \cdot a^{-1} \equiv 1 \pmod{p}.
    $$
    Uniqueness follows from cancellation: if $ax \equiv 1 \equiv ay \pmod{p}$, then $ax \equiv ay$ and $a \not\equiv 0 \pmod{p}$ give $x \equiv y \pmod{p}$.
}

\EnvId{agpNrS7Gk-plIDAp2Xc0d-inverses-add-like-fractions}
\Theorem{Inverses Add Like Fractions}{
    Let $b, d \not\equiv 0 \pmod{p}$. Then for any $a, c$, we have
    $$\frac{a}{b} + \frac{c}{d} \equiv a \cdot b^{-1} + c \cdot d^{-1} \equiv (ad + bc) \cdot (bd)^{-1} \equiv \frac{ad + bc}{bd} \pmod{p},$$
    just like normal fractions.
}{
    The first and last congruences are just fraction-bar notation, so we prove the middle one. Since $p$ is a prime and $b, d \not\equiv 0 \pmod{p}$, we also have $bd \not\equiv 0 \pmod{p}$, and hence $(bd)^{-1}$ exists by the previous theorem.

    Multiplying by $bd$ and using $b^{-1}b \equiv 1$ and $d^{-1}d \equiv 1$ gives
    $$
    \bigl(a b^{-1} + c d^{-1}\bigr) \cdot bd \equiv ad + bc \pmod{p},
    $$
    and multiplying by $(bd)^{-1}$ gives $a b^{-1} + c d^{-1} \equiv (ad + bc) \cdot (bd)^{-1} \pmod{p}$.
}

\EnvId{4FiBDjLCxOcmWRSNkUVRs-inverses-multiply-like-fractions}
\Theorem{Inverses Multiply Like Fractions}{
    Let $b, d \not\equiv 0 \pmod{p}$. Then for any $a, c$, we have
    $$\frac{a}{b} \cdot \frac{c}{d} \equiv (a \cdot b^{-1}) \cdot (c \cdot d^{-1}) \equiv (ac) \cdot (bd)^{-1} \equiv \frac{ac}{bd} \pmod{p},$$
    just like normal fractions.
}{
    The outer congruences are just fraction-bar notation, so we prove the middle one. As in the previous theorem, $bd \not\equiv 0 \pmod{p}$, so $(bd)^{-1}$ exists. Moreover,
    $$
    (bd) \cdot \bigl(b^{-1} d^{-1}\bigr) \equiv \bigl(b b^{-1}\bigr)\bigl(d d^{-1}\bigr) \equiv 1 \pmod{p},
    $$
    so by uniqueness of the inverse, $(bd)^{-1} \equiv b^{-1} d^{-1} \pmod{p}$. From this,
    $$
    \bigl(a b^{-1}\bigr) \cdot \bigl(c d^{-1}\bigr) \equiv (ac) \cdot \bigl(b^{-1} d^{-1}\bigr) \equiv (ac) \cdot (bd)^{-1} \pmod{p}.
    $$
}

\sec Problems

The problems below all benefit from the theory above, working with fractions modulo~$p$ as if they were ordinary numbers. They are ordered roughly by difficulty.

\EnvId{xwufjUXqxotXL6cgvQ4gW-wilson-factorial-congruence}
\Problem{0}{Wilson}{
    Let $p$ be a prime. Prove that
    $$(p-1)! \equiv -1 \pmod{p}.$$
}{
    In the product $1 \cdot 2 \cdots (p-1)$, try to pair factors with their inverse. Which numbers do not have a pair?
}{
    The numbers without the pair are those satisfying $k \equiv 1/k \mod p$. Find them all.
}{
    For $p = 2$ we have $(p-1)! = 1 \equiv -1 \pmod{2}$, so assume from now on that $p$ is odd.

    Every $k \in \{1, 2, \ldots, p-1\}$ is nonzero modulo $p$, so it has an inverse $1/k$, which is nonzero as well and hence congruent to exactly one element of $\{1, 2, \ldots, p-1\}$. The assignment $k \mapsto 1/k$ is therefore a map from this set to itself and is its own inverse. It thus splits the set into pairs $\{k, 1/k\}$ with $k \not\equiv 1/k$, whose product is congruent to $1$ by the definition of the inverse, and into fixed points, that is, those $k$ with
    $$
    k \equiv \frac{1}{k} \pmod{p}.
    $$
    Multiplying by $k$, we see that this condition is equivalent to $k^2 \equiv 1 \pmod{p}$, that is, $p \mid (k-1)(k+1)$. Since $p$ is a prime, it divides one of the factors, hence $k \equiv 1$ or $k \equiv -1 \pmod{p}$. In $\{1, 2, \ldots, p-1\}$ these correspond exactly to $k = 1$ and $k = p-1$, which are distinct for odd $p$.

    The remaining $p-3$ numbers therefore split into $(p-3)/2$ pairs, each with product congruent to $1$, and so
    $$
    (p-1)! \equiv 1 \cdot (p-1) \cdot 1^{(p-3)/2} \equiv -1 \pmod{p}.
    $$
}

\EnvId{oZgQ3gdX2sSdnfW0T-I1N-prime-dividing-sum-of-squares}
\Problem{0}{}{
    Let $p$ be a prime. Prove that there exist integers $a, b$, neither divisible by $p$, with $p \mid a^2 + b^2$ if and only if $p = 2$ or $p \equiv 1 \pmod{4}$.
}{
    The condition $p \mid a^2 + b^2$ rewrites as $(a/b)^2 \equiv -1 \pmod{p}$. So the question becomes: for which primes is $-1$ a square modulo~$p$?
}{
    For one direction, raise the relation $(a/b)^2 \equiv -1$ to the $(p-1)/2$-th power. The left side becomes $(a/b)^{p-1}$.
}{
    Apply Fermat's Little Theorem to the left side. What does the resulting equation force on $p \mod 4$?
}{
    For the converse with $p \equiv 1 \pmod{4}$, you need an explicit square root of $-1$ modulo~$p$. A natural source: Wilson's theorem, $(p-1)! \equiv -1 \pmod{p}$.
}{
    In the product $(p-1)!$, pair $k$ with $p - k \equiv -k$. What does the result look like for $p \equiv 1 \pmod{4}$?
}{
    For $p = 2$, the pair $a = b = 1$ works. From now on let $p$ be an odd prime and put $m = \frac{p-1}{2}$.

    We first rewrite the condition in the language of fractions modulo $p$. If $p \nmid b$, then $p \mid a^2 + b^2$ means $a^2 \equiv -b^2 \pmod{p}$, and dividing by $b^2$ gives
    $$
    \bigl(\tfrac{a}{b}\bigr)^2 \equiv -1 \pmod{p}.
    $$
    Conversely, if some integer $x$ satisfies $x^2 \equiv -1 \pmod{p}$, then the pair $a = x$, $b = 1$ works, since $x^2 \equiv -1$ immediately gives $p \nmid x$. So the pair we are looking for exists if and only if $-1$ is a square modulo $p$.

    Suppose such a pair exists and put $x = \frac{a}{b}$, so that $x^2 \equiv -1 \pmod{p}$. Raising to the $m$-th power gives
    $$
    x^{p-1} = \bigl(x^2\bigr)^m \equiv (-1)^m \pmod{p}.
    $$
    Fermat's Little Theorem applied to $a$ and to $b$ gives $x^{p-1} = \frac{a^{p-1}}{b^{p-1}} \equiv 1 \pmod{p}$, hence $(-1)^m \equiv 1 \pmod{p}$. Were $m$ odd, we would get $-1 \equiv 1 \pmod{p}$, that is, $p \mid 2$, which fails for odd $p$. So $m$ is even, and that is exactly the condition $p \equiv 1 \pmod{4}$.

    Conversely, let $p \equiv 1 \pmod{4}$, so that $m$ is even. We need an explicit square root of $-1$ modulo~$p$, and Wilson's theorem $(p-1)! \equiv -1 \pmod{p}$ from the previous problem supplies it. The pairs $\{k, p-k\}$ for $k = 1, 2, \ldots, m$ partition $1, 2, \ldots, p-1$, and $p - k \equiv -k \pmod{p}$, so
    $$
    (p-1)! = \prod_{k=1}^{m} k(p-k) \equiv (-1)^m (m!)^2 \pmod{p}.
    $$
    Since $m$ is even, $(-1)^m = 1$, and Wilson's theorem thus gives $(m!)^2 \equiv -1 \pmod{p}$. Hence the pair $a = m!$, $b = 1$ works: $a$ is a product of numbers smaller than $p$, so $p \nmid a$, and $a^2 + b^2 \equiv -1 + 1 \equiv 0 \pmod{p}$.
}

\EnvId{iQkNpjFlIqY_8YqY2fOcG-divisibility-of-a2-ab-b2}
\Problem{0}{}{
    Let $p \equiv 2 \pmod{3}$ be a prime. Show that if $p \mid a^2 + ab + b^2$, then $p \mid a$ and $p \mid b$.
}{
    Note that $a^3 - b^3 = (a - b)(a^2 + ab + b^2)$, so $p \mid a^3 - b^3$. Assume for contradiction that $p \nmid b$. This is similar to the previous problem.
}{
    Raise to the $(p-2)/3$-th power, which is a non-negative integer since $p \equiv 2 \pmod{3}$. The left side becomes $(a/b)^{p-2}$.
}{
    Multiply both sides by $a/b$ to get $(a/b)^{p-1} \equiv a/b \pmod{p}$. Apply Fermat's Little Theorem to conclude $a \equiv b \pmod{p}$. Substitute this back.
}{
    Assume for contradiction that $p \nmid b$. Then $a$ is not divisible by $p$ either: if $p \mid a$, then $p \mid a^2 + ab + b^2$ immediately gives $p \mid b^2$, hence $p \mid b$.

    The key is the factorization
    $$
    a^3 - b^3 = (a - b)(a^2 + ab + b^2),
    $$
    which gives $p \mid a^3 - b^3$, that is, $a^3 \equiv b^3 \pmod{p}$. Since $p \nmid b$, we may divide this congruence by $b^3$, and for the residue $x = a/b$ we get
    $$
    x^3 \equiv 1 \pmod{p}.
    $$

    This is where the assumption $p \equiv 2 \pmod{3}$ comes in: the number $(p-2)/3$ is a non-negative integer, so we may raise the last congruence to that power and obtain
    $$
    x^{p-2} \equiv \bigl(x^3\bigr)^{(p-2)/3} \equiv 1 \pmod{p}.
    $$
    Multiplying by $x$ turns this relation into $x^{p-1} \equiv x \pmod{p}$. Moreover, $p \nmid a$, so $x \not\equiv 0 \pmod{p}$, and Fermat's Little Theorem gives $x^{p-1} \equiv 1 \pmod{p}$. Comparing the two congruences, we get
    $$
    \frac{a}{b} = x \equiv 1 \pmod{p}, \qquad\text{hence}\qquad a \equiv b \pmod{p}.
    $$

    Substituting this back into the assumption gives
    $$
    0 \equiv a^2 + ab + b^2 \equiv 3b^2 \pmod{p}.
    $$
    Since $p \nmid b$, we must have $p \mid 3$, that is, $p = 3$. This contradicts the assumption $p \equiv 2 \pmod{3}$.

    Therefore $p \mid b$. Then $a^2 = (a^2 + ab + b^2) - b(a + b)$ is divisible by $p$, and since $p$ is a prime, $p \mid a$ as well.
}

\EnvId{5dOxVJIKK4L9pkK7lsJb8-binomial-coefficient-p-minus-1}
\Problem{0}{}{
    Let $p$ be a prime and let $0 \leq k \leq p-1$ be an integer. Prove that
    $$\binom{p-1}{k} \equiv (-1)^k \pmod{p}.$$
}{
    Expand the binomial coefficient as $\binom{p-1}{k} = \frac{(p-1)(p-2) \cdots (p-k)}{k!}$. 
}{
    Reduce each factor in the numerator modulo $p$: $p - j \equiv -j \pmod{p}$. Aren't we left with something familiar?
}{
    We expand the binomial coefficient as a fraction and read it modulo~$p$. Since $0 \leq k \leq p-1$, all the factors $1, 2, \ldots, k$ are smaller than $p$, and therefore $k! \not\equiv 0 \pmod{p}$; by the theorem on the existence of the inverse, $k!$ is invertible modulo~$p$ and the expression
    $$
    \binom{p-1}{k} = \frac{(p-1)(p-2)\cdots(p-k)}{k!}
    $$
    makes sense modulo~$p$ as a fraction.

    The numerator has exactly $k$ factors of the form $p - j$ for $j = 1, 2, \ldots, k$, and each of them satisfies $p - j \equiv -j \pmod{p}$. Pulling a minus sign out of each factor gives
    $$
    (p-1)(p-2)\cdots(p-k) \equiv (-1)(-2)\cdots(-k) = (-1)^k \cdot k! \pmod{p},
    $$
    so the numerator is, up to sign, the same factorial as the denominator. Dividing, we get
    $$
    \binom{p-1}{k} \equiv \frac{(-1)^k \cdot k!}{k!} \equiv (-1)^k \pmod{p}.
    $$
    For $k = 0$ both products above are empty, hence equal to $1$, and the claim reads $1 \equiv 1 \pmod{p}$.
}

\EnvId{QHLcqpD-GBM-mlx7Gk8Oa-reciprocal-sum-numerator}
\Problem{0}{}{
    Let $p \geq 3$ be a prime. Show that if
    $$\frac{1}{1} + \frac{1}{2} + \frac{1}{3} + \cdots + \frac{1}{p-1} = \frac{m}{n},$$
    then $p \mid m$.
}{
    Looking at this modulo~$p$, we have the sum of inverses of all possible non-zero values mod $p$.
}{
    Realize that this sum is actually the sum of all values from $1$ to $p-1$, not necessarily in this order.
}{
    We first put the sum over the common denominator $N = (p-1)!$. All the numbers $N/k$ for $k = 1, 2, \ldots, p-1$ are integers, so the integer $M = \sum_{k=1}^{p-1} N/k$ satisfies
    $$
    \frac{1}{1} + \frac{1}{2} + \cdots + \frac{1}{p-1} = \frac{M}{N} = \frac{m}{n},
    $$
    hence $mN = nM$. The number $N$ is the product of $1, 2, \ldots, p-1$, none of which is divisible by $p$, so $p \nmid N$. Therefore, if we prove $p \mid M$, we get $p \mid nM = mN$, and since $N$ is coprime to $p$, $p \mid m$ as well, which is the claim. Note that we assume nothing about the form of the fraction $m/n$.

    It remains to compute $M$ modulo $p$. For every $k$ we have $k \cdot \frac{N}{k} = N$, so modulo $p$ the integer $N/k$ is exactly $N \cdot k^{-1}$, that is, the fraction $N/k$ in the sense of inverses. Therefore
    $$
    M \equiv N \sum_{k=1}^{p-1} \frac{1}{k} \pmod{p}.
    $$

    The sum on the right is the sum of the inverses of all non-zero residues modulo $p$. The map $k \mapsto k^{-1}$ is a permutation of these residues: from $k \cdot k^{-1} \equiv 1 \pmod{p}$ we see that the inverse of $k^{-1}$ is again $k$, so it is an involution and hence a bijection. The numbers $1^{-1}, 2^{-1}, \ldots, (p-1)^{-1}$ are therefore, in some order, exactly the numbers $1, 2, \ldots, p-1$, and
    $$
    \sum_{k=1}^{p-1} \frac{1}{k} \equiv 1 + 2 + \cdots + (p-1) = \frac{p(p-1)}{2} \pmod{p}.
    $$
    Since $p \geq 3$ is odd, $(p-1)/2$ is an integer and the right-hand side is a multiple of $p$. We get $\sum_{k=1}^{p-1} 1/k \equiv 0$, hence $M \equiv N \cdot 0 \equiv 0 \pmod{p}$, so $p \mid M$, and by the opening argument $p \mid m$.
}

\EnvId{4ctqjmewZ7WcyqJtpQrAp-imo-2005-coprime-to-sequence}
\Problem{0}{IMO 2005}{
    Determine all positive integers $n$ that are coprime to every term of the sequence
    $$a_m = 2^m + 3^m + 6^m - 1, \quad m \geq 1.$$
}{
    Show that for every prime $p$, some $a_m$ is divisible by $p$. Then the only valid $n$ is $n = 1$.
}{
    Small primes are easy: $a_1 = 10$ handles $p = 2, 5$ and $a_2 = 48$ handles $p = 3$. Focus on $p \geq 5$ and try to choose $m$ depending on $p$ so that the four terms collapse.
}{
    For $p \geq 5$, take $m = p - 2$. Then $2^m \equiv 1/2$, $3^m \equiv 1/3$, $6^m \equiv 1/6 \pmod{p}$. Compute $a_{p-2} \pmod{p}$.
}{
    \Answer{The only such $n$ is $n = 1$.}

    We show that every prime $p$ divides some term of the sequence. That settles it: if $n > 1$, it has a prime divisor $p$, that divisor divides some $a_m$, and hence $n$ and $a_m$ are not coprime.

    We handle the small primes directly. We have $a_1 = 2 + 3 + 6 - 1 = 10$, which is divisible by $2$ and by $5$, and $a_2 = 4 + 9 + 36 - 1 = 48$, which is divisible by $3$.

    So let $p \geq 5$. The numbers $2$, $3$, $6$ are not divisible by $p$, so modulo $p$ we can work with them as fractions. The key to the choice of $m$ is the identity
    $$
    \frac{1}{2} + \frac{1}{3} + \frac{1}{6} = 1.
    $$
    If the powers $2^m$, $3^m$, $6^m$ turned into exactly $1/2$, $1/3$, $1/6$, the whole term $a_m$ would come out zero. The choice $m = p - 2$, which is positive since $p \geq 5$, does exactly that. Indeed, by Fermat's Little Theorem
    $$
    2 \cdot 2^{p-2} = 2^{p-1} \equiv 1 \pmod{p},
    $$
    so $2^{p-2}$ is the inverse of $2$, that is, $2^{p-2} \equiv \frac{1}{2} \pmod{p}$. Similarly $3^{p-2} \equiv \frac{1}{3}$ and $6^{p-2} \equiv \frac{1}{6} \pmod{p}$. Adding these inverses as fractions gives
    $$
    \gather
    a_{p-2} \equiv \frac{1}{2} + \frac{1}{3} + \frac{1}{6} - 1 \\
    \equiv \frac{3 + 2 + 1}{6} - 1 \equiv 1 - 1 \equiv 0 \pmod{p}.
    \endgather
    $$

    Every prime therefore divides some term of the sequence, and $n = 1$ is the only candidate left. And indeed, the number $1$ is coprime to every integer, hence to every $a_m$.
}

\EnvId{h6JPweRuH3hf9vtoEksW0-product-of-k2-k-1}
\Problem{0}{Poland}{
    Let $p > 3$ be a prime with $p \equiv 2 \pmod{3}$, and let $a_k = k^2 + k + 1$ for $k = 1, 2, \ldots, p-1$. Prove that
    $$a_1 a_2 \cdots a_{p-1} \equiv 3 \pmod{p}.$$
}{
    For $k \neq 1$, factor $k^2 + k + 1 = \frac{k^3 - 1}{k - 1}$. Isolate the $k = 1$ term, which contributes~$3$.
}{
    The product becomes $3 \cdot \frac{\prod_{k=2}^{p-1}(k^3 - 1)}{(p-2)!}$. We need to calculate the numerator and denominator mod~$p$ separately.
}{
    It would be great if the product of all the numbers $k^3-1$ was actually something we could calculate. 
}{
    We isolate the term $a_1 = 3$. For $k \geq 2$ we have $k - 1 \not\equiv 0 \pmod{p}$, so we may divide by this number modulo $p$, and the identity $k^3 - 1 = (k-1)(k^2+k+1)$ gives
    $$
    a_k \equiv \frac{k^3 - 1}{k - 1} \pmod{p}.
    $$
    Over all $k = 2, 3, \ldots, p-1$ the denominators multiply to $1 \cdot 2 \cdots (p-2) = (p-2)!$, so
    $$
    a_1 a_2 \cdots a_{p-1} \equiv 3 \cdot \frac{\prod_{k=2}^{p-1} (k^3 - 1)}{(p-2)!} \pmod{p}.
    $$

    The numerator looks intractable, but the product of \textit{all} numbers of the form $x - 1$ over the non-zero residues $x$ would be easy to compute. It therefore suffices to know which residues the cubes run through, and this is exactly where the assumption on~$p$ modulo $3$ comes in.

    We show that the map $x \mapsto x^3$ is a bijection on the non-zero residues modulo $p$. Since $p \equiv 2 \pmod{3}$, the number $t = \frac{p-2}{3}$ is a positive integer. Suppose $x^3 \equiv y^3 \pmod{p}$ for non-zero residues $x$, $y$. Raising to the $t$-th power gives $x^{p-2} \equiv y^{p-2} \pmod{p}$, which by Fermat's Little Theorem reads
    $$
    \frac{1}{x} \equiv \frac{1}{y} \pmod{p},
    $$
    hence $x \equiv y \pmod{p}$. An injective map of a finite set into itself is a bijection.

    The numbers $1^3, 2^3, \ldots, (p-1)^3$ are therefore a permutation of $1, 2, \ldots, p-1$ modulo $p$. So the only $k$ with $k^3 \equiv 1 \pmod{p}$ is $k = 1$, and for $k = 2, 3, \ldots, p-1$ the values $k^3$ run through exactly the residues $2, 3, \ldots, p-1$ in some order. We get
    $$
    \prod_{k=2}^{p-1} (k^3 - 1) \equiv \prod_{j=2}^{p-1} (j - 1) = (p-2)! \pmod{p}.
    $$
    The number $(p-2)!$ is not divisible by $p$, so it cancels in the fraction and we are left with
    $$
    a_1 a_2 \cdots a_{p-1} \equiv 3 \pmod{p}.
    $$
}

\EnvId{XmHaSAnFcLBJMQwLRRJu6-power-sum-over-half-range}
\Problem{0}{Singapore 2012}{
    Let $p$ be an odd prime. Prove that
    $$\sum_{i=1}^{(p-1)/2} i^{p-2} \equiv \frac{2 - 2^p}{p} \pmod{p}.$$
}{
    Since $i^{p-2} \equiv 1/i \pmod{p}$, the left-hand side is $\sum_{i=1}^{(p-1)/2} 1/i$. Now work on the right-hand side to get the hint.
}{
    Expand $2^p = (1+1)^p = \sum_{k=0}^{p} \binom{p}{k}$, so $\frac{2^p - 2}{p} = \frac{1}{p} \sum_{k=1}^{p-1} \binom{p}{k}$.
}{
    Use $\binom{p}{k} = \frac{p}{k} \binom{p-1}{k-1}$ to rewrite each summand: $\frac{2^p - 2}{p} = \sum_{k=1}^{p-1} \frac{\binom{p-1}{k-1}}{k}$.
}{
    Apply $\binom{p-1}{k-1} \equiv (-1)^{k-1} \pmod{p}$ from a~previous problem to simplify.
}{
    The right-hand side reduces to $-\sum_{k=1}^{p-1} (-1)^{k-1}/k \pmod{p}$. Split this into even and odd $k$, and use $\sum_{k=1}^{p-1} 1/k \equiv 0 \pmod{p}$ from a~previous problem.
}{
    First note that the right-hand side makes sense: by Fermat's Little Theorem $2^p \equiv 2 \pmod{p}$, so $p \mid 2^p - 2$ and $\frac{2-2^p}{p}$ is an integer.

    The left-hand side is immediate. For $1 \leq i \leq p-1$ we have $i \not\equiv 0 \pmod{p}$, so $i^{p-2} \equiv i^{p-1} \cdot \frac{1}{i} \equiv \frac{1}{i} \pmod{p}$. Therefore
    $$
    \sum_{i=1}^{(p-1)/2} i^{p-2} \equiv \sum_{i=1}^{(p-1)/2} \frac{1}{i} \pmod{p}. \eqno(1)
    $$
    All the work is on the right-hand side, and the goal is to extract that same sum from it.

    The binomial theorem gives $2^p = (1+1)^p = \sum_{k=0}^{p} \binom{p}{k}$, so after removing the two end terms
    $$
    \frac{2^p-2}{p} = \frac{1}{p} \sum_{k=1}^{p-1} \binom{p}{k}.
    $$
    We can divide the sum by $p$ term by term thanks to the identity $\binom{p}{k} = \frac{p}{k} \binom{p-1}{k-1}$ (both sides equal $\frac{p!}{k!(p-k)!}$), after which $p$ cancels:
    $$
    \frac{2^p-2}{p} = \sum_{k=1}^{p-1} \frac{\binom{p-1}{k-1}}{k}. \eqno(2)
    $$
    This is an equality of rational numbers, but every denominator on the right-hand side satisfies $1 \leq k \leq p-1$, hence is not divisible by $p$. So we may read (2) as a congruence modulo~$p$, where $\frac{1}{k}$ means the inverse of $k$.

    By an earlier problem, $\binom{p-1}{k-1} \equiv (-1)^{k-1} \pmod{p}$, so multiplying (2) by $-1$ turns it into
    $$
    \frac{2-2^p}{p} \equiv \sum_{k=1}^{p-1} \frac{(-1)^k}{k} \pmod{p}.
    $$

    It remains to split this alternating sum by the parity of $k$. The number $p-1$ is even, so the even indices are exactly $k = 2j$ for $j = 1, \ldots, \frac{p-1}{2}$; put
    $$
    E = \sum_{j=1}^{(p-1)/2} \frac{1}{2j} \equiv \frac{1}{2} \sum_{j=1}^{(p-1)/2} \frac{1}{j} \pmod{p}
    $$
    and $S = \sum_{k=1}^{p-1} \frac{1}{k}$. The sum over odd $k$ is then $S - E$, hence
    $$
    \sum_{k=1}^{p-1} \frac{(-1)^k}{k} = E - (S - E) = 2E - S.
    $$
    By an earlier problem, $S \equiv 0 \pmod{p}$, so
    $$
    \frac{2-2^p}{p} \equiv 2E \equiv \sum_{j=1}^{(p-1)/2} \frac{1}{j} \pmod{p},
    $$
    which together with (1) is exactly the congruence to be proved.
}

\EnvId{FfwCijxt7LgTb4ro14tSW-binomial-2p-choose-p}
\Problem{0}{}{
    Show that for every prime $p$,
    $$\binom{2p}{p} \equiv 2 \pmod{p^2}.$$
}{
    Write $\binom{2p}{p} = \frac{(2p)(2p-1)\cdots(p+1)}{p!}$. Pull out the factor $2p$ in the numerator and the $p$ in $p! = p \cdot (p-1)!$ to get $\binom{2p}{p} = \frac{2 \prod_{k=1}^{p-1}(p+k)}{(p-1)!}$.
}{
    Factor each numerator term as $p + k = k\bigl(1 + \frac{p}{k}\bigr)$, so $\prod_{k=1}^{p-1}(p+k) = (p-1)! \cdot \prod_{k=1}^{p-1}\bigl(1 + \frac{p}{k}\bigr)$.
}{
    Expand $\prod\bigl(1 + \frac{p}{k}\bigr)$ modulo $p^2$: only the constant and linear-in-$p$ terms survive, giving $1 + p \sum_{k=1}^{p-1} 1/k$.
}{
    By P4, $\sum_{k=1}^{p-1} 1/k \equiv 0 \pmod{p}$, so $p \sum 1/k \equiv 0 \pmod{p^2}$.
}{
    For $p = 2$ we check the claim directly: $\binom{4}{2} = 6 = 4 + 2 \equiv 2 \pmod{4}$. So from now on let $p \geq 3$.

    Expand the binomial coefficient and pull the multiples of $p$ out of it:
    $$
    \binom{2p}{p} = \frac{(2p)(2p-1)\cdots(p+1)}{p!}.
    $$
    In the numerator the only multiple of $p$ is the first factor $2p$, the remaining factors being $p + k$ for $k = 1, 2, \ldots, p-1$; the denominator is $p! = p \cdot (p-1)!$. Cancelling one $p$ gives the integer identity
    $$
    \binom{2p}{p} \cdot (p-1)! = 2 \prod_{k=1}^{p-1}(p+k). \eqno(1)
    $$

    We now compute the right-hand side modulo~$p^2$. For $1 \leq k \leq p-1$ the number $k$ is coprime to $p^2$, so it has an inverse modulo~$p^2$ as well, and we write $\frac1k$ for this inverse; the proofs of both theorems on computing with fractions use nothing but the invertibility of the denominator, so they go through verbatim modulo~$p^2$. From $k \cdot \frac{p}{k} = p$ we get $p + k \equiv k\bigl(1 + \frac{p}{k}\bigr) \pmod{p^2}$, hence
    $$
    \prod_{k=1}^{p-1}(p+k) \equiv (p-1)! \prod_{k=1}^{p-1}\bigl(1 + \tfrac{p}{k}\bigr) \pmod{p^2}.
    $$

    Expand the last product. The term corresponding to a subset $S \subseteq \{1, \ldots, p-1\}$ is $p^{|S|} \prod_{k \in S} \frac1k$, which for $|S| \geq 2$ is divisible by $p^2$ and hence zero modulo~$p^2$. Only the $1$ and the terms linear in~$p$ survive:
    $$
    \prod_{k=1}^{p-1}\bigl(1 + \tfrac{p}{k}\bigr) \equiv 1 + p \sum_{k=1}^{p-1} \frac1k \pmod{p^2}.
    $$

    It therefore suffices to show that $p \sum_{k=1}^{p-1} \frac1k \equiv 0 \pmod{p^2}$, that is, that the sum $\sum_{k=1}^{p-1} \frac1k$ itself is divisible by~$p$. Only information modulo~$p$ matters here, and the inverse of $k$ modulo~$p^2$ becomes the inverse of $k$ modulo~$p$ once reduced modulo~$p$. So modulo~$p$ our sum is an ordinary sum of inverses modulo~$p$, and that sum is zero for $p \geq 3$: we proved it in the problem on the numerator of the sum $\frac11 + \frac12 + \cdots + \frac1{p-1}$. This is exactly where $p = 2$ fails, because for $p = 2$ this sum equals $1$.

    The product in the last relation is thus $1$ modulo~$p^2$, and (1) turns into
    $$
    \binom{2p}{p} \cdot (p-1)! \equiv 2 \cdot (p-1)! \pmod{p^2}.
    $$
    The number $(p-1)!$ is coprime to $p^2$, so we may cancel it, which gives $\binom{2p}{p} \equiv 2 \pmod{p^2}$.
}

\EnvId{IyalyrEDv2e9srzrTKrZS-reciprocal-power-sum-numerator}
\Problem{0}{}{
    Let $p \geq 5$ be a prime and let $e$ be a positive integer with $(p-1) \nmid e$. Show that if
    $$\frac{1}{1^e} + \frac{1}{2^e} + \frac{1}{3^e} + \cdots + \frac{1}{(p-1)^e} = \frac{m}{n},$$
    then $p \mid m$.
}{
    The claim is $\sum_{k=1}^{p-1} \frac{1}{k^e} \equiv 0 \pmod{p}$. The trick is to look for an $a$ such that multiplying every $k$ by $a$ shuffles the terms but rescales the whole sum.
}{
    Use the assumption $(p-1) \nmid e$ to find $a$ with $a^e \not\equiv 1 \pmod{p}$. The trick is to consider a primitive root.
}{
    First we turn the claim about the numerator into a congruence. Set $N = \bigl((p-1)!\bigr)^e$. For every $k \in \{1, 2, \ldots, p-1\}$, $k^e$ divides $N$, so putting the sum over the common denominator $N$ gives
    $$
    \frac{1}{1^e} + \frac{1}{2^e} + \cdots + \frac{1}{(p-1)^e} = \frac{M}{N},
    $$
    where $M = \sum_{k=1}^{p-1} \frac{N}{k^e}$ is an integer. The equality $\frac{m}{n} = \frac{M}{N}$ gives $mN = nM$, and $N$ is a product of numbers smaller than $p$, hence $p \nmid N$. To prove $p \mid m$ it therefore suffices to prove $p \mid M$.

    The integer $\frac{N}{k^e}$ is congruent modulo~$p$ to the product $N \cdot \frac{1}{k^e}$, where $\frac{1}{k^e}$ is the inverse of $k^e$ modulo~$p$: since $k^e \not\equiv 0 \pmod{p}$, this congruence is equivalent by cancellation to $N \equiv N$. So if we denote
    $$
    S = \frac{1}{1^e} + \frac{1}{2^e} + \cdots + \frac{1}{(p-1)^e},
    $$
    where the fractions are now read as inverses modulo~$p$, then $M \equiv N \cdot S \pmod{p}$. Since $p \nmid N$, it suffices to prove $S \equiv 0 \pmod{p}$.

    We look for an $a \not\equiv 0 \pmod{p}$ for which replacing every $k$ by $ak$ shuffles the terms but rescales the sum as a whole. The shuffling works for any such $a$: the residues $a \cdot 1, a \cdot 2, \ldots, a(p-1)$ are a permutation of $1, 2, \ldots, p-1$ modulo~$p$, and since the value of $\frac{1}{k^e}$ modulo~$p$ depends only on the residue of $k$,
    $$
    S \equiv \sum_{k=1}^{p-1} \frac{1}{(ak)^e} \pmod{p}.
    $$
    The rescaling comes from the theorem on multiplying fractions applied to $(ak)^e = a^e k^e$:
    $$
    \sum_{k=1}^{p-1} \frac{1}{a^e k^e}
    \equiv \frac{1}{a^e} \sum_{k=1}^{p-1} \frac{1}{k^e}
    \equiv \frac{S}{a^e} \pmod{p}.
    $$
    Putting the two together, we get $S \equiv \frac{S}{a^e}$, so multiplying by $a^e$ gives
    $$
    (a^e - 1) \cdot S \equiv 0 \pmod{p}.
    $$
    Hence if we find an $a$ with $a^e \not\equiv 1 \pmod{p}$, the factor $a^e - 1$ is non-zero and cancellation yields $S \equiv 0 \pmod{p}$.

    This is exactly what the assumption $(p-1) \nmid e$ is for. Let $g$ be a primitive root modulo~$p$, that is, an element of order $p-1$, and set $a = g$. Write $e = q(p-1) + r$ with $0 \leq r < p-1$; the assumption gives $r \neq 0$. By Fermat's Little Theorem,
    $$
    g^e = \bigl(g^{p-1}\bigr)^q \cdot g^r \equiv g^r \pmod{p},
    $$
    so $g^e \equiv 1 \pmod{p}$ would force $g^r \equiv 1 \pmod{p}$ with $0 < r < p-1$, contradicting the fact that the order of $g$ is $p-1$. Hence $g^e \not\equiv 1 \pmod{p}$ and the proof is complete.
}

\EnvId{5GgUO9sDTBS1pWxire2HE-wolstenholme-reciprocal-sum}
\Problem{0}{Wolstenholme}{
    Let $p \geq 5$ be a prime. Show that if
    $$\frac{1}{1} + \frac{1}{2} + \frac{1}{3} + \cdots + \frac{1}{p-1} = \frac{m}{n},$$
    then $p^2 \mid m$.
}{
    Pair $k$ with $p - k$.
}{
    This gives $\sum_{k=1}^{p-1} 1/k = p \sum_{k=1}^{(p-1)/2} 1/(k(p-k))$. We need the right-hand factor to be divisible by~$p$.
}{
    Modulo $p$, $1/(k(p-k)) \equiv -1/k^2$. So it suffices that $\sum_{k=1}^{(p-1)/2} 1/k^2 \equiv 0 \pmod{p}$.
}{
    Note $(p-k)^2 \equiv k^2 \pmod{p}$, so $\sum_{k=1}^{p-1} 1/k^2 = 2 \sum_{k=1}^{(p-1)/2} 1/k^2$. Can we now use a previous problem?
}{
    Denote
    $$
    S = \frac{1}{1} + \frac{1}{2} + \cdots + \frac{1}{p-1}, \qquad D = (p-1)!.
    $$
    Throughout, $S$ is understood as an ordinary fraction, not as a residue modulo~$p$; congruences will only be applied to auxiliary integers. Note at the outset that $p \nmid D$.

    Adding the terms $1/k$ and $1/(p-k)$ produces a $p$ in the numerator:
    $$
    \frac{1}{k} + \frac{1}{p-k} = \frac{p}{k(p-k)}.
    $$
    Since $p$ is odd, the pairs $\{k, p-k\}$ for $k = 1, 2, \ldots, \frac{p-1}{2}$ cover each of $1, 2, \ldots, p-1$ exactly once, hence
    $$
    S = p \sum_{k=1}^{(p-1)/2} \frac{1}{k(p-k)}.
    $$

    One factor of $p$ comes for free, and it suffices to show that the remaining sum contributes another one. This is a statement about a fraction, so let us put it over the common denominator $D$. For $k \leq \frac{p-1}{2}$, the numbers $k$ and $p-k$ are two distinct elements of $\{1, 2, \ldots, p-1\}$, so $k(p-k) \mid D$ and the number
    $$
    A = \sum_{k=1}^{(p-1)/2} \frac{D}{k(p-k)}
    $$
    is an integer, with $S = \frac{pA}{D}$. We want to prove $p \mid A$.

    Only here do we start working with residues. Each summand $X = \frac{D}{k(p-k)}$ satisfies $X \cdot k(p-k) = D$, so modulo~$p$ we have $X \equiv D \cdot \bigl(k(p-k)\bigr)^{-1}$. Moreover $k(p-k) \equiv -k^2 \pmod{p}$, and therefore
    $$
    A \equiv -D \sum_{k=1}^{(p-1)/2} \frac{1}{k^2} \pmod{p}.
    $$
    Since $p \nmid D$, it remains to prove that
    $$
    \sum_{k=1}^{(p-1)/2} \frac{1}{k^2} \equiv 0 \pmod{p},
    $$
    where this is a sum of residues, that is, of inverses in the sense of the theory above.

    We extend this half sum to the full one, for which we already know the claim. From $p - k \equiv -k$ we get $(p-k)^2 \equiv k^2$, hence also $\frac{1}{(p-k)^2} \equiv \frac{1}{k^2}$, so the second half of the terms is a copy of the first:
    $$
    \sum_{k=1}^{p-1} \frac{1}{k^2} \equiv 2 \sum_{k=1}^{(p-1)/2} \frac{1}{k^2} \pmod{p}.
    $$
    For $p \geq 5$ we have $p - 1 \geq 4$, so $(p-1) \nmid 2$, and by the previous problem with $e = 2$ the left-hand side is zero modulo~$p$. The number $2$ is coprime to $p$, so the half-range sum is zero as well, and indeed $p \mid A$.

    Write $A = pB$ with $B$ an integer, so $S = \frac{p^2 B}{D}$. Multiplying out $\frac{m}{n} = \frac{p^2 B}{D}$ gives $mD = p^2 Bn$, so $p^2 \mid mD$. Since $p \nmid D$, we get $p^2 \mid m$.
}

\bye
